\section{Analysis \& Discussion}
\label{sec:discussion}

This section interprets the results from Section~\ref{sec:results},
identifies the key patterns that emerge from the benchmark, and
acknowledges the limitations and threats to validity that constrain
the conclusions we can draw.

%% ----------------------------------------------------------------
\subsection{Key Findings}
\label{sec:key-findings}

\paragraph{LLM scanners represent a qualitative shift, not an incremental gain.}
The most striking result in Table~\ref{tab:main-results} is not that
LLM-based scanners outperform traditional SAST tools---it is the
\emph{magnitude} of the gap.  The best LLM scanner (Sonnet~4, F2\,=\,53.8)
achieves nearly 3$\times$ the F2 score of the best traditional tool
(Snyk, F2\,=\,18.2), and across the board LLM scanners deliver 2--4$\times$
higher F2 scores than their SAST counterparts.  The per-CWE breakdown
reveals where this advantage originates: on vulnerability classes that
require semantic understanding of data flow---SQL injection (91\% vs.\
32\% recall), command injection (78\% vs.\ 24\%), and insecure
deserialization---LLMs dramatically outperform pattern-based tools.
Traditional SAST tools remain competitive only on issues that reduce to
syntactic patterns (e.g., hardcoded secrets), but even on those
categories their overall recall remains low.  These results suggest that
LLM-based analysis is not merely a refinement of existing SAST
methodology but a fundamentally different approach that can reason about
code semantics in ways that rule-based tools cannot.

\paragraph{Agentic tool use is highly cost-effective.}
Comparing agentic and single-turn modes of the same underlying model
provides a controlled test of the value of tool use in vulnerability
detection.  Agentic Haiku achieves an F2 of 39.4 versus 27.0 for
single-turn Haiku---a 46\% improvement---while incurring only 6\% higher
cost.  The ability to read files, grep for patterns, and iteratively
explore a codebase allows the model to gather context that a single
prompt cannot provide.  This is consistent with the broader finding that
retrieval-augmented and tool-augmented LLM workflows outperform
prompt-only setups in code understanding tasks.  From a practical
standpoint, the marginal cost of agentic execution is small relative to
the gain in detection quality, making it the dominant strategy for
LLM-based scanning.

\paragraph{Bigger models do not always win; reliability matters.}
One might expect a monotonic relationship between model capability
(and cost) and benchmark performance.  The data tell a more nuanced
story.  Opus~4 achieves the highest micro F2 score on the repositories
it completes (62.4), but its 36\% timeout rate (completing only 19 of 26
repos) drags its strict micro F2 down to 49.9---below Sonnet~4's 53.8,
which reflects full coverage.  A similar pattern emerges with
Minimax, which times out on 31\% of repos.  Conversely,
Kimi~k2 and Minimax~M1 punch above their weight on
cost-efficiency metrics, delivering competitive F2 scores at
substantially lower per-repo cost.  The implication for practitioners is
that scanner selection should optimize for \emph{reliability-adjusted}
performance rather than peak capability: a model that reliably completes
every scan is more valuable than one that occasionally excels but
frequently fails to return results.

\paragraph{The precision--recall tradeoff is stark and F2 captures it correctly.}
Grok~4.20 Reasoning exemplifies the conservative end of the spectrum:
its precision of 0.927 is the highest in the benchmark, but its recall
of 0.263 means it misses nearly three-quarters of true vulnerabilities.
Most LLM scanners cluster in the 0.6--0.8 precision range with
substantially higher recall, which the F2 metric---designed to weight
recall twice as heavily as precision---correctly rewards.  Among
traditional tools, SonarQube illustrates the opposite failure mode:
respectable precision (0.611) paired with negligible recall (0.065),
yielding an F2 of just 7.2.  The F2 metric aligns well with
practitioner priorities in security scanning, where missing a real
vulnerability (false negative) is generally more costly than
investigating a false alarm (false positive).  Benchmarks that report
only precision or only detection rate would paint a misleading picture
of scanner utility.

%% ----------------------------------------------------------------
\subsection{Limitations}
\label{sec:limitations}

We identify seven limitations that scope the conclusions of this work:

\begin{itemize}
  \item \textbf{Python only.}
    Version~1.0 of \textsc{RealVuln} covers only Python repositories.
    Scanner performance may differ substantially on statically typed
    languages (Java, Go) or memory-unsafe languages (C/C++) where
    vulnerability classes such as buffer overflows are prevalent.
    Extension to additional languages is planned for v2
    (Section~\ref{sec:living-benchmark}).

  \item \textbf{Type~1 repositories only.}
    All 26 repos are intentionally vulnerable applications with
    artificially high vulnerability density.  Production codebases have
    sparser vulnerabilities embedded in larger, more complex code, which
    may alter both precision and recall characteristics.  Type~2--5
    targets (Section~\ref{sec:target-types}) will address this gap.

  \item \textbf{Vendor conflict of interest.}
    Kolega, which leads the benchmark at F2\,=\,66.5, is also the
    organization publishing this work.  We mitigate this concern by
    open-sourcing the entire benchmark: ground truth labels, scoring
    code, scanner outputs, and the dashboard.  Any researcher can
    independently verify our results or challenge the ground truth.

  \item \textbf{Ground truth subjectivity.}
    Hand-labeling vulnerabilities inherently involves judgment calls,
    particularly for borderline cases such as informational findings or
    context-dependent misconfigurations.  We publish the labeling
    evidence and rationale for every ground truth entry and invite
    community challenges via the project's issue tracker.

  \item \textbf{LLM non-determinism.}
    LLM-based scanners produce variable output across runs.  We report
    multi-run statistics (mean, standard deviation) where available, but
    cost constraints prevented multi-run evaluation for every model
    configuration.  Single-run results should be interpreted with
    appropriate uncertainty.

  \item \textbf{Default scanner configurations.}
    All scanners are evaluated using default or vendor-recommended
    settings.  Tuned configurations---custom rules for Semgrep, adjusted
    quality profiles for SonarQube, or modified system prompts for
    LLM scanners---could yield different results.  We report the exact
    configuration used for each scanner to enable reproduction.

  \item \textbf{Timeouts and incomplete scans.}
    Opus~4 timed out on approximately 36\% of repositories and
    Minimax~M1 on approximately 31\%.  These incomplete scans reduce
    effective coverage and are reflected in the strict micro F2 metric,
    which treats unscanned repos as zero.  The gap between lenient and
    strict scores quantifies the reliability penalty.
\end{itemize}

%% ----------------------------------------------------------------
\subsection{Threats to Validity}
\label{sec:threats}

\paragraph{Data contamination.}
The 26 target repositories are publicly available on GitHub, and it is
likely that some or all of them appeared in the training data of the LLM
scanners we evaluate.  If a model has memorized known vulnerabilities in
these repos, its recall could be inflated relative to performance on
unseen code.  We mitigate this threat in two ways.  First, the 120
false-positive traps in the ground truth test whether scanners
\emph{discriminate} between vulnerable and safe code; memorizing that a
file contains vulnerabilities does not help a scanner avoid flagging a
non-vulnerable code path in the same file.  Second, the v2 roadmap
(Section~\ref{sec:living-benchmark}) includes repositories published
after model training cutoffs, which will provide a contamination-free
evaluation set.

\paragraph{Selection bias.}
Twenty-six repositories, while substantial for a hand-labeled benchmark,
do not exhaustively represent the Python vulnerability landscape.
Certain CWE families (e.g., cryptographic weaknesses, race conditions)
are underrepresented in the current corpus.  The living-benchmark
roadmap addresses this through planned expansion to additional repos
and target types, aiming for broader CWE and application-domain
coverage in future versions.
